%%% -*- coding: utf-8 -*-
% !TeX root = main-slides.tex

\lecture{Networking}{unity-networking}
\section{Networking}

\showtitlepage{Netcode for GameObjects}

\begin{frame}{One Program, Several Roles}
  \framelabel{slide:roles}
  \begin{itemize}
    \item Netcode for GameObjects replicates GameObjects across the peers of a
      session. One peer is the server; a \emph{host} is server and client in one
      process, which is how SEE is usually run.
    \item The same script runs on every peer. A networked script therefore has
      to ask which role it is in, and SEE does so throughout:
      \texttt{IsServer}, \texttt{IsHost}, \texttt{IsClient},
      \texttt{IsOwner}.
    \item Only a GameObject carrying a \texttt{NetworkObject} component takes
      part at all, and only its \texttt{NetworkBehaviour} components can
      replicate anything.
    \item There are two mechanisms and no more: state that replicates itself
      (\texttt{NetworkVariable}) and calls that travel to other peers (RPCs).
  \end{itemize}
\end{frame}

\begin{frame}[fragile]{NetworkBehaviour}
  \begin{lstlisting}[basicstyle=\scriptsize\ttfamily]
// Assets/SEE/Net/ActionNetwork.cs
public class ActionNetwork : NetworkBehaviour
{
    public override void OnNetworkSpawn() { ... }
}

// Assets/SEE/Game/Avatars/PlayerName.cs
private void Start()   // not OnNetworkSpawn: for a dynamically
{                      // spawned object, Start runs after it
    if (IsLocalPlayer) { ... }
}
  \end{lstlisting}
  \vspace{0.5ex}

  \begin{itemize}
    \item A \texttt{NetworkBehaviour} is a \texttt{MonoBehaviour} with the
      network lifecycle added. \texttt{OnNetworkSpawn} and
      \texttt{OnNetworkDespawn} bracket the time the object is replicated, which
      is not the same span as \texttt{Awake} to \texttt{OnDestroy}.
    \item The order bites: for a dynamically spawned object \texttt{Start} runs
      \emph{after} \texttt{OnNetworkSpawn}, which is why \texttt{PlayerName}
      waits until \texttt{Start} for data the server fills in while spawning.
  \end{itemize}
\end{frame}

\begin{frame}[fragile]{NetworkVariable}
  \begin{lstlisting}[basicstyle=\scriptsize\ttfamily]
// Assets/SEE/Game/Avatars/AvatarHandAnimationsSync.cs
private readonly NetworkVariable<Vector3> leftHandPosition
    = new(writePerm: NetworkVariableWritePermission.Owner);

// the owner writes, every other peer only reads
if (IsOwner)
{
    CaptureFromHandsAnimator();   // assigns .Value
}
else { ... }                      // reads .Value
  \end{lstlisting}
  \vspace{0.5ex}

  \begin{itemize}
    \item A \texttt{NetworkVariable<T>} replicates its \texttt{Value} to every
      peer whenever it changes. Reading is always allowed, writing needs the
      permission given at construction: the server by default, or the owner.
    \item SEE hands the hand poses to the owner, because they come from that
      user's webcam, and the peers that do not own the avatar only read.
    \item \texttt{T} has to be something Netcode can serialize without
      allocating, such as \texttt{Vector3}, \texttt{Quaternion},
      \texttt{float} or a \texttt{FixedString}. Plain \texttt{string} is not
      among them.
  \end{itemize}
\end{frame}

\begin{frame}[fragile]{Remote Procedure Calls}
  \begin{lstlisting}[basicstyle=\scriptsize\ttfamily]
// Assets/SEE/Game/Avatars/PlayerName.cs
[Rpc(SendTo.Server)]
private void RequestPlayerNameFromServerRpc(ulong networkObjectId,
                                            RpcParams rpcParams = default)
{
    ulong senderClientId = rpcParams.Receive.SenderClientId;
    string playerName = PlayerNameMap.GetPlayerName(networkObjectId);
    SendPlayerNameToClientRpc(networkObjectId, playerName,
        RpcTarget.Single(senderClientId, RpcTargetUse.Temp));
}
[Rpc(SendTo.SpecifiedInParams)]
void SendPlayerNameToClientRpc(ulong networkObjectId,
                               string playerName, RpcParams _)
  \end{lstlisting}
  \vspace{0.5ex}

  \begin{itemize}
    \item \texttt{[Rpc]} marks a method that is called on one peer and runs on
      others, chosen by \texttt{SendTo}. The name must end in \texttt{Rpc}; the
      compiler insists on it and generates the transport.
    \item The optional \texttt{RpcParams} carry who sent the call, and let the
      server answer that one client alone. Here a client asks the server for a
      name and the server sends it back, a round trip in two methods.
  \end{itemize}
\end{frame}

\begin{frame}[fragile]{Ownership and Authority}
  \framelabel{slide:authority}
  \begin{lstlisting}[basicstyle=\scriptsize\ttfamily]
// Assets/SEE/Net/ClientNetworkTransform.cs
public class ClientNetworkTransform : NetworkTransform
{
    // Owner client only. This imposes state to the server.
    protected override bool OnIsServerAuthoritative()
    {
        return false;
    }
}
  \end{lstlisting}
  \vspace{0.5ex}

  \begin{itemize}
    \item Ownership belongs to the \texttt{NetworkObject}: exactly one peer owns
      it, the server unless stated otherwise, while a player's avatar is owned
      by that player's client.
    \item Authority is the question of who may change what. A
      \texttt{NetworkTransform} is server-authoritative by default, so a client
      moving its own avatar would see it corrected back.
    \item SEE overrides that, and its own comment names the price: the owning
      client imposes the transform on the server, which means trusting clients,
      so nothing security-sensitive may ride on it.
  \end{itemize}
\end{frame}

\begin{exerciseframe}{Which Mechanism?}
  \begin{enumerate}
    \item An avatar's hand pose has to be visible to everyone and changes
      continuously while the hand moves. Which of the two mechanisms, and who
      is allowed to write?
    \item A player's name is a \texttt{string}. Why can it not be a
      \texttt{NetworkVariable<string>}, and what does SEE do instead?
    \item You drag your own avatar and it snaps back to where it was. What is
      wrong, and which method does SEE override to prevent it?
  \end{enumerate}
\end{exerciseframe}

\begin{solutionframe}{Which Mechanism?}
  \begin{enumerate}
    \item A \texttt{NetworkVariable}, written by the owner: continuous state
      that everyone reads is what it is for, and an RPC per frame per joint
      would be traffic for nothing.
    \item Netcode replicates a \texttt{NetworkVariable} without allocating, so
      \texttt{T} cannot be a reference type. RPC parameters may be strings,
      hence \texttt{PlayerName} requests the name from the server instead.
    \item The transform is server-authoritative, so the server corrects what
      the client did. \texttt{ClientNetworkTransform} overrides
      \texttt{OnIsServerAuthoritative} to \texttt{false}
      (\slideref{slide:authority}).
  \end{enumerate}
\end{solutionframe}
